\documentclass[11pt]{article}
\usepackage{amsmath,amssymb,amsthm,mathtools,booktabs}
\usepackage[margin=1in]{geometry}
\usepackage[hidelinks]{hyperref}
\newcommand{\W}{\mathcal W}
\newcommand{\C}{\mathbb C}
\newcommand{\R}{\mathbb R}
\newcommand{\diag}{\operatorname{diag}}
\title{Complete-Weil positivity for the cubic Gaussian source and certified movable-pole matrices}
\author{}
\date{20 September 2026}
\begin{document}
\maketitle

\section{Original source, full zero sum and the human inputs}
Retain the completed zeta function and the actual Gaussian source
\[
 \xi(s)=\Phi(s)\zeta(s),\qquad
 \Phi(s)=\tfrac12s(s-1)\pi^{-s/2}\Gamma(s/2),\qquad
 \phi_*(x)=(4\pi^2x^4-6\pi x^2)e^{-\pi x^2}.                 \tag{WP1}
\]
For $D=-x\,d/dx$, $P(s)=\sum_{j=0}^3p_js^j$ and
$f_P=P(D)\phi_*$, integration by parts gives
\[
 F_P(s)=\int_0^\infty f_P(x)x^{s-1}\,dx=P(s)\Phi(s),
                  \qquad \Re s>-2.                            \tag{WP2}
\]
Indeed the two terms of the Mellin integral of $\phi_*$ are
$2\pi^{-s/2}\Gamma(s/2+2)$ and
$-3\pi^{-s/2}\Gamma(s/2+1)$; the Gamma recursion gives WP1.
Every integration by parts has zero boundary term on the stated
half-plane, since $D^j\phi_*(x)=O(x^2)$ at zero and is a polynomial
times a Gaussian at infinity. The even extension of every $f_P$ is
Schwartz, vanishes at zero, and has integral zero, since $\Phi(1)=0$.

All sums below run over the actual nontrivial zeros $\rho$ of $\zeta$,
with their actual multiplicities $m_\rho$. Define
\[
 \W(F,G)=\sum_\rho m_\rho\,
       \overline{F(1-\bar\rho)}G(\rho).                         \tag{WP3}
\]
The functional equation and conjugation preserve this multiset.
The estimates below prove absolute convergence for the sources
used here, so reflection and conjugation may be applied without
an ordering convention. At a critical-line zero each summand of
$\W(F,F)$ is $m_\rho|F(\rho)|^2$.

The finite-height input is Platt--Trudgian's Theorem~1
\cite{PlattTrudgian}: every zero up to the stated height
$3\,000\,175\,332\,800$ is on the critical line. We use only
\[
                         T_0=3\cdot10^{12}.                    \tag{WP4}
\]
For the unverified tail use Hasanalizade--Shen--Wong's
Corollary~1.2 \cite{HSW}, in its stated original version:
\[
 \left|N(T)-\frac{T}{2\pi}\log\frac{T}{2\pi e}\right|
 \le .1038\log T+.2573\log\log T+9.3675\quad(T\ge e).          \tag{WP5}
\]
$N(T)$ counts positive-height zeros with multiplicity. In particular
$N(T)\le T\log T$ for $T\ge100$: the right-hand error in WP5 is
less than $\log T+10\le T$, while
$1/(2\pi)+1/\log T<1$. These inequalities hold at $100$ and
persist because $T-\log T-10$ increases and $1/\log T$ decreases.
The proof uses no hypothesis on the location of zeros above WP4.

\section{Four certified zeros and the Gamma estimate}
The script \texttt{certify\_four\_critical\_zeros.py}, using Arb
ball arithmetic \cite{Johansson} at $384$ bits, encloses the full
WP1 expression at $s=1/2+it$. It proves the following strict real
signs; the displayed decimal intervals are outward coarse
enclosures of the stored balls, not replacement calculations:
\[
\begin{array}{c|c}
t&\xi(1/2+it)\\\hline
14&(0.0002012,0.0002014)\\
15&(-0.0007058,-0.0007056)\\
21&(-3.978\cdot10^{-7},-3.977\cdot10^{-7})\\
22&(7.716\cdot10^{-6},7.717\cdot10^{-6})
\end{array}                                                     \tag{WP6}
\]
The functional equation and conjugation show that this function
of real $t$ is real. Continuity and the sign changes therefore
produce zeros $\gamma_1\in(14,15)$ and $\gamma_2\in(21,22)$.
Together with their negatives these give four distinct zeros
$\rho_j=1/2+i\gamma_j$, where here the index runs over
$\gamma_1,\gamma_2,-\gamma_1,-\gamma_2$. No simplicity is needed.
Every pair of these four heights has distance greater than $6$.
The complete balls, script hash and library version are retained
in \texttt{FOUR\_CRITICAL\_ZERO\_CERTIFICATES.json}.

For clarity the Gamma constants are proved explicitly from
the complex Stirling remainder in Askey--Roy, DLMF 5.11.1 and
5.11(ii) \cite{Gamma}. In the open right half-plane,
\[
 \log\Gamma z=(z-\tfrac12)\log z-z+\tfrac12\log(2\pi)+R(z),
 \quad |R(z)|\le\frac{\sec^2(\arg z/2)}{12|z|}
                   \le\frac1{6|z|}.                           \tag{WP7}
\]
The first inequality is the remainder after omitting all
Bernoulli terms. Put $z=\beta/2+it/2$ with $0\le\beta\le1$,
$t\ge10$, and $w=z+1=a+ib$, so $1\le a\le3/2$, $b\ge5$.
Using $b\arg w\ge b\pi/2-a$, WP7 and $\Gamma(z)=\Gamma(w)/z$,
\[
 \begin{split}
 |\Gamma(z)|
 &\le \sqrt{2\pi}\,b^{\beta/2-1/2}
       (1+.09)^{1/2}e^{1/30}e^{-\pi b/2}
   <3e^{-\pi t/4},\\
 |\Phi(\beta+it)|&\le5(1+t^2)e^{-\pi t/4}.
 \end{split}                                                   \tag{WP8}
\]
Here $|w|^{a-1/2}=b^{a-1/2}(1+a^2/b^2)^{(a-1/2)/2}$,
$|z|\ge b$, $a^2/b^2\le .09$, and $\beta/2-1/2\le0$.
Also $|s(s-1)|\le1+t^2$ and $\pi^{-\beta/2}\le1$.
The same estimates hold for negative $t$ by conjugation.

For any of the four critical heights, apply WP7 directly to
$z=1/4+i\gamma/2$, so $7<|z|<12$. It gives
\[
 \log|\Gamma z|^2\ge
 \log(2\pi)-\tfrac12\log12-11\pi-\tfrac12-\tfrac1{21}>-36.
\]
Since $|s(s-1)|^2\ge1$ at those heights,
\[
             |\Phi(\rho_j)|^2>e^{-40}\quad(1\le j\le4).
                                                                  \tag{WP9}
\]

\section{The full cubic Gaussian inequality, with all high zeros retained}
Let $V_{j\ell}=\rho_j^\ell$ for $0\le\ell\le3$ and let
$H=V^*\diag(|\Phi(\rho_j)|^2)V$ be the contribution of one
copy of each of the four zeros. It is positive definite, and
\[
 \det H=\prod_{j=1}^4|\Phi(\rho_j)|^2
                   \prod_{i<j}|\rho_j-\rho_i|^2>e^{-160},
 \qquad \operatorname{tr}H<400\cdot485^5<e^{40}.              \tag{WP10}
\]
The trace estimate follows from WP8 and
$\sum_{\ell=0}^3|\rho_j|^{2\ell}\le4(1+\gamma_j^2)^3$,
discarding only the factor $e^{-\pi|\gamma_j|/2}\le1$ in this
upper bound. Hence
\[
                         \lambda_{\min}(H)>e^{-280}.           \tag{WP11}
\]
Indeed the product of its three largest eigenvalues is at most
$(\operatorname{tr}H)^3$.

For a unit coefficient vector $p$, Cauchy--Schwarz gives
$|P(\beta+it)|\le2(1+t^2)^{3/2}$ uniformly in the strip.
The absolute value of each unknown zero summand in WP3 is at
most
\[
                  u(t)=100(1+t^2)^5e^{-ct},\qquad c=\pi/2.
                                                                  \tag{WP12}
\]
For $t\ge T_0$, $u$ decreases and
$0\le-u'(t)\le c u(t)\le3200c\,t^{10}e^{-ct}$.
Stieltjes integration by parts, WP5 and $\log t\le t$ show
\[
 \begin{split}
 2\sum_{\Im\rho>T_0}m_\rho u(\Im\rho)
 &\le2\int_{T_0}^{\infty}N(t)(-u'(t))\,dt\\
 &\le6400e^{-cT_0}\sum_{j=0}^{12}
             \frac{12!}{(12-j)!}\frac{T_0^{12-j}}{c^j}
 <e^{-T_0}.
 \end{split}                                                   \tag{WP13}
\]
For the last inequality the sum is at most
$T_0^{12}\sum_{j=0}^{12}(12/(cT_0))^j<2T_0^{12}$.
Thus its logarithm including $6400$ is less than
$10+12\log T_0<370$, whereas $(c-1)T_0>10^{12}$.
The factor $2$ includes negative heights; the boundary term in
integration by parts is nonpositive and was discarded in an
upper bound. WP12--13 also prove absolute convergence.

All other zeros below WP4 contribute nonnegative terms, by
Platt--Trudgian. Combining WP11 and WP13 and then scaling $p$
proves the complete, unconditional inequality
\[
 \boxed{\W(\Phi P,\Phi P)\ge e^{-300}\sum_{j=0}^3|p_j|^2.}
                                                                  \tag{WP14}
\]
The margin follows from $e^{-280}-e^{-T_0}>e^{-300}$.
The unknown high zeros have been estimated, not placed on the
critical line. WP14 is a statement about this entire original
cubic Gaussian source space, not only four sampled vectors.

\section{Exact translation and the coefficient-space receiver}
The original target coordinate is $S'=c_*+iy$, whereas the
zeta coordinate on the same physical line is $s=1/2+iy$.
Thus $S'=s+c_*-1/2$. Retain this exact affine map, and define
the actual coefficient translation $C_{c_*}$ by
$p(s+c_*-1/2)=\sum_i(C_{c_*}p)_i s^i$:
\[
 (C_{c_*})_{ij}=\begin{cases}\binom ji(c_*-1/2)^{j-i},&i\le j,\\
                  0,&i>j,\end{cases}\qquad
 \det C_{c_*}=1,\quad C_{c_*}^{-1}=C_{1-c_*}.                \tag{WP15}
\]
The entry formula is the binomial theorem; composing the two
translations proves the inverse. For the original conductor
coordinate map $y=\Psi q$ of FC26--32 \cite{FC} into the coefficient vector of the
cubic $p$, the original lift is exactly $p(D+c_*-1/2)\phi_*$.
Explicitly, $\Psi=YK^{-1}$, where $K$ is the retained collision
coordinate matrix and the $a$th column of $Y$ is the coefficient
vector of the original polynomial
\[
 y_a(S')=\sum_{u,w=0}^{8}a_{uw}
          (S'+\beta_{8,u,w})^v\ell_a(S'+\beta_{8,u,w}).
\]
Here $\ell_a(S)=\prod_{b\ne a}(S-\rho_b)/(\rho_a-\rho_b)$
and $v$ is the actual order of the conductor symbol, exactly as
in FC17 and FC26; its periods and original coefficients remain
in $a_{uw}$ and $\beta_{8,u,w}$.
Writing $\mathsf W_\Phi$ for the coefficient Gram in WP14,
its pullback is therefore
\[
 \boxed{\mathsf W_{\rm original}
  =\Psi^*C_{c_*}^*\mathsf W_\Phi C_{c_*}\Psi
  \succeq e^{-300}\Psi^*C_{c_*}^*C_{c_*}\Psi.}              \tag{WP16}
\]
This identity follows by applying WP2 to each original
coefficient column; it retains every entry of $\Psi$ and the
centre. The arithmetic native metric and its full minimum are
not replaced by the Euclidean coefficient norm: WP16 specifies
the actual form to be transported to those coordinates.
For the separate eight-state conductor, its original target map
is substituted into this displayed pullback by the Fable
receiver cited in the accompanying source index.

\section{A rational cubic test and its unknown-tail bound}
Retain the pole $\omega=2$, its $h=3$, and the four functions
\[
 U(s)=\frac{s+1}{2-s},\qquad
 \varphi_j(s)=\frac{\sqrt3}{2-s}U(s)^j\quad(0\le j\le3).
                                                                  \tag{WP17}
\]
The complete movable-pole proof MP in the companion source \cite{MP}
establishes that their full-Weil Gram is the Hermitian Toeplitz
matrix $T_3(2)$; equivalently this matrix is defined directly by
WP3 applied to these four functions. The present positivity
proof uses WP3 and therefore does not depend on a truncated
zero expansion or on a Taylor approximation to its entries.

At a critical height $t$, the weight of one row is
$w_t=3/(9/4+t^2)>1/200$ for $14<|t|<22$, and $|U(1/2+it)|=1$.
For two of our four heights,
\[
 |U(1/2+it)-U(1/2+iu)|
 =\frac{3|t-u|}{\sqrt{9/4+t^2}\sqrt{9/4+u^2}}
 >\frac{18}{487}>\frac1{30}.                                \tag{WP18}
\]
Consequently the four-row Vandermonde Gram $H_R$ satisfies
\[
 \det H_R> b_*:=\frac1{200^4 30^{12}},\qquad
 \operatorname{tr}H_R\le\frac{16\cdot3}{9/4+14^2}<1,
 \qquad\lambda_{\min}(H_R)>b_*.                              \tag{WP19}
\]
On $0\le\Re s\le1$, $|\Im s|=t\ge10$, one has
$|U(s)|^2\le1+3/t^2<1.02^2$ and
$|U'(s)|\le3/t^2$. Differentiation of WP17 gives, for each
$0\le j\le3$,
\[
 |\varphi_j'(s)|
 \le\frac{\sqrt3}{t^2}
       \left(1.02^3+\frac9t\,1.02^2\right)<\frac4{t^2}.
\]
Thus for $F=\sum_{j=0}^3p_j\varphi_j$ with $\|p\|_2=1$,
\[
 |F'(s)|\le8/t^2,\qquad
 |F(\rho)-F(1-\bar\rho)|\le8/t^2.                           \tag{WP20}
\]
The second bound integrates on the horizontal segment of length
$|2\Re\rho-1|\le1$ inside the same strip. A noncritical reflected
pair with equal multiplicity $m$ contributes
\[
 2m\Re(\overline{F(1-\bar\rho)}F(\rho))
 =m\left(|F(\rho)|^2+|F(1-\bar\rho)|^2
        -|F(\rho)-F(1-\bar\rho)|^2\right)
 \ge-64m/t^4.                                                 \tag{WP21}
\]
Each positive-height reflected pair accounts for $2m$ in $N$;
the same bound at negative heights therefore gives a total
possible negative contribution of at most
\[
 \begin{split}
 64\sum_{\Im\rho>T_0}\frac{m_\rho}{(\Im\rho)^4}
 &\le\frac{256}{T_0^3}
                    \left(\frac{\log T_0}{3}+\frac19\right)
 <10^{-33}<\frac{b_*}{2}.
 \end{split}                                                   \tag{WP22}
\]
Here Stieltjes integration by parts gives an upper bound
$4\int_{T_0}^\infty N(t)t^{-5}\,dt$, followed by WP5. The last
comparison follows already from $\log T_0<30$ and
$200^4 30^{12}<10^{27}$. Absolute convergence of WP3 also follows
from $F(s)=O(1/|\Im s|)$ and $N(t)=O(t\log t)$.
All verified zeros outside the selected four contribute
nonnegatively. Hence
\[
                      \boxed{T_3(2)\succeq\tfrac12 b_* I.}    \tag{WP23}
\]
This directly bounds all four rational coefficients, including
their arbitrary relative phases.

\section{New exact-ball certificates through degree eight}
For a pole $\omega=2+i\tau$ put
\[
 \begin{gathered}
 h=3,\quad M_\omega(z)=\frac{\omega+(1-\bar\omega)z}{1+z},
 \quad L=\xi'/\xi,\\
 L(M_\omega(z))=L(\omega)+\sum_{n\ge1}c_nz^n,
 \quad c_0=2\Re L(\omega),\quad c_{-n}=\bar c_n.
 \end{gathered}                                                \tag{WP24}
\]
The companion MP proof identifies $T_d(\omega)=(c_{j-i})_{0\le i,j\le d}$
with the complete WP3 Gram of
$\sqrt h\,(\omega-s)^{-1}[(s+\bar\omega-1)/(\omega-s)]^j$.
At these poles $\Re\omega>1$, so $\xi(\omega)\ne0$ and the
Taylor-series division is valid. The checked script
\texttt{certify\_weil\_matrices.py} computes WP24 by the full
WP1 product using Arb, with $768$ bits and series cap $12$.
If $X(z)=\xi(M_\omega(z))$, it computes
\[
          L(M_\omega(z))=\frac{X'(z)}{X(z)M_\omega'(z)}.
                                                                  \tag{WP25}
\]
Since $M_\omega'(0)=-h\ne0$, series division is again valid.
The series cap retains every coefficient through degree eight
after differentiation. All constants are exact integers or
Arb-enclosed transcendental constants and values; no zero list
or finite prime sum is substituted into WP25.

For each $\tau\in\{0,14,21,100,1000\}$, the script constructs
all leading principal matrices of sizes $1$ through $9$ and
encloses each determinant. Every determinant has strictly
positive real lower endpoint; its imaginary ball contains
zero, as required by the exact Hermitian symmetry. Sylvester's
criterion therefore proves
\[
 \boxed{T_d(2+i\tau)\succ0\quad
       (0\le d\le8,\ \tau\in\{0,14,21,100,1000\}).}           \tag{WP26}
\]
This is $45$ positive determinant certificates. The script
also encloses and proves positive each of the $40$ margins
$c_0-|c_n|$ for $1\le n\le8$. The complete ball coefficients,
all minors, all margins, precision, version and script hash are
in \texttt{ARITHMETIC\_BALL\_CERTIFICATES.json}. For orientation
only, the determinants of the largest matrices are approximately
\[
\begin{array}{c|r}
\tau&\det T_8(2+i\tau)\\\hline
0&2.87069033117\cdot10^{-79}\\
14&2.04193998712\cdot10^{-52}\\
21&1.04788150392\cdot10^{-37}\\
100&4.21694999273\cdot10^{-13}\\
1000&96.1156305520
\end{array}                                                     \tag{WP27}
\]
The strict ball comparisons, rather than these rounded
numbers, are the certificates used in WP26.

WP14, WP23 and WP26 establish full-form positive statements on
the displayed finite source spaces, including the unknown
zero tails. The movable-pole degree remains unrestricted in
the exact MP criterion. No conclusion about all degrees, or
about RH, is inferred from these finite certificates.

\begin{thebibliography}{9}
\bibitem{MP} Split-Zero programme,
\emph{Exact movable-pole Weil coefficients and finite arithmetic
certificates}, MP1--38, complete accompanying source
\texttt{MOVABLE\_POLE\_WEIL\_DERIVATION.tex}.
\bibitem{FC} Split-Zero programme,
\emph{Fable to the original conductor}, FC17 and FC26--32,
complete source \texttt{FABLE\_TO\_ORIGINAL\_CONDUCTOR.tex},
published edition:
\url{https://github.com/KokunoYumeto/zeta-function-research-reader/tree/f8b73284d2e669e981bf2c7609e99f4d2b3a50be/workbenches/splitzero-tandem/continuations/20260920-marked-observation-illustrated}.
\bibitem{PlattTrudgian} D.~Platt and T.~Trudgian,
\emph{The Riemann hypothesis is true up to $3\cdot10^{12}$},
arXiv:2004.09765v1, Theorem~1; Bull. London Math. Soc.,
doi:10.1112/blms.12460.
\url{https://arxiv.org/abs/2004.09765v1}.
Original author TeX read completely; canonical disk source
\texttt{PUBUNIT-8FD792C15B5C6472FBDED93E}.
\bibitem{HSW} E.~Hasanalizade, Q.~Shen and P.-J.~Wong,
\emph{Counting zeros of the Riemann zeta function},
arXiv:2107.06506v1, Corollary~1.2 and its displayed explicit
bound. \url{https://arxiv.org/abs/2107.06506v1}.
The original TeX introduction and this corollary were read;
this citation imports that published zero-counting theorem.
\bibitem{Gamma} R.~A.~Askey and R.~Roy,
\emph{Gamma Function}, NIST Digital Library of Mathematical
Functions, Chapter~5, formula 5.11.1 and section 5.11(ii),
release 1.2.8. Original formula TeX retained.
\url{https://dlmf.nist.gov/5.11.E1};
\url{https://dlmf.nist.gov/5.11.ii}.
\bibitem{Johansson} F.~Johansson,
\emph{Arb: efficient arbitrary-precision midpoint-radius
interval arithmetic}, IEEE Transactions on Computers 66
(2017), 1281--1292, doi:10.1109/TC.2017.2690633.
Original author source: \url{https://arxiv.org/abs/1611.02831v1},
sections 1, 2, 5 and the power-series/matrix parts of section 6.
Arb is used through python-flint 0.9.0; the precision and
strict enclosure tests are recorded in both supplied scripts.
\bibitem{Lagarias} J.~C.~Lagarias,
\emph{Li Coefficients for Automorphic $L$-Functions},
arXiv:math/0404394v4, Lemmas 2.1 and 4.1; Ann. Inst. Fourier
57 (2007), 1689--1740.
\url{https://arxiv.org/abs/math/0404394v4}.
The original genus-one product and convergence proof are the
human primary inputs to the complete movable-pole companion.
\end{thebibliography}
\end{document}
